% Propietario conservado: /Users/ruben/Documents/New project/output/REVISION_ARGUMENTAL_APERTURAS_CIERRES_20260915/IX/ES/sections/ch_profundidad_constantes.tex
\section{Profondeur coinductive de la quadrirelation}
\label{ix:sec:profundidad-cuatrirrelacion-ch}

Les lecteurs régionaux et la clôture dodécaphasique publient, respectivement, \(\pi,\varphi,e\) et \(\alpha\) à partir de l'état enrichi commun. Nous réservons \(N\) à la profondeur opératorielle et \(n\) à la longueur décimale : les deux indices ne sont pas identifiés et aucun coût uniforme de lecture n'est présupposé. Pour chaque réalisation exacte \(\chi\), sa partie entière est conservée séparément et l'on note \(p_{\chi,n}\) le préfixe de longueur \(n\) de sa partie fractionnaire, dans le développement décimal canonique non égal à neuf à partir d'un certain rang. On définit
\[
 \rho^{\rm dec}_{n+1,n}(d_1\ldots d_nd_{n+1})=d_1\ldots d_n,
 \qquad p_{\chi,0}:=\varnothing.
\]
L’unicité de \(p_{\pi,n},p_{\varphi,n},p_{e,n}\) provient de l’unicité du
caractère exact déjà engendré et de son cylindre décimal canonique. Celle de
\(p_{\alpha,n}\) provient de la clôture entière dodécaphasique, de son identité
des queues et de la compatibilité projective prouvée dans le
théorème~\ref{ix:thm:alpha-salida-dodecafasica-completa}. La résolvante analytique
d’ordre neuf apporte une caractérisation finie supplémentaire, tandis que le
lecteur analytique complet calcule les coefficients restants au moyen de
\(b_{10+3n}=\lambda/729^n\), prouve la convergence uniforme et construit
des intervalles rationnels inclus dans les cylindres entiers. Il identifie ainsi
cette même section à toute profondeur par le
théorème~\ref{ix:prop:alpha-dos-limites}. Le jet fini n’est utilisé ni comme
générateur ni comme substitut du prolongement coinductif.

\begin{theorem}[Compatibilité projective de \(\pi,\varphi,e,\alpha\) à toute profondeur]
\label{ix:thm:profundidad-cuatrirrelacion-ch}
L’action coinductive de l’état commun détermine
\[
 \begin{gathered}
  \forall\chi\in\{\pi,\varphi,e,\alpha\}\;
  \forall n\ge1\quad\exists!\,p_{\chi,n},\\
  \rho^{\rm dec}_{n+1,n}(p_{\chi,n+1})=p_{\chi,n},
  \qquad
  p_{\chi,\infty}:=(p_{\chi,n})_{n\ge0}
  \in\varprojlim_{n\ge0}\{0,\ldots,9\}^{n}.
 \end{gathered}
\tag{D1}
\label{ix:eq:profundidad-cuatrirrelacion-ch}
\]
En particulier,
\[
 (p_{\pi,\infty},p_{\varphi,\infty},p_{e,\infty},p_{\alpha,\infty})
\tag{D2}
\label{ix:eq:cuatrirrelacion-infinita-ch}
\]
est une publication coinductive corrélée de l’unique état enrichi.
\end{theorem}

\begin{proof}
Le résultat~\ref{ix:gen:prefijos-regionales} fixe les réalisations exactes
des trois lecteurs régionaux. Le
théorème~\ref{ix:thm:alpha-salida-dodecafasica-completa} construit la section
complète \(\alpha_{\mathrm{HMT}}\) et prouve la compatibilité de tous ses
préfixes. Le développement décimal canonique de chacune de ces
quatre réalisations possède un unique préfixe de chaque longueur ; tronquer le
préfixe de longueur \(n+1\) donne celui de longueur \(n\). Appliquer ces lecteurs
à l’état commun donne la famille corrélée énoncée. L’existence et la
compatibilité de cette publication exacte n’identifient pas \(n\) à \(N\)
et n’ajoutent pas de borne de coût aux procédés de lecture déjà établis.
\end{proof}

L'ordre constructif interne conserve les lecteurs de \(\pi,\varphi,e\), l'état signé, le sceau dodécaphasique entier \(K\) et la roue dodécaphasique qui publie \(\alpha\) ; aucune constante conventionnelle n'est promue au rang de générateur. Une exécution à mille, dix mille ou un million de chiffres vérifie l'instance computationnelle finie correspondante de \eqref{ix:eq:profundidad-cuatrirrelacion-ch} ; elle ne constitue pas une borne du théorème. Le théorème est la famille compatible pour tout \(n\) et la section infinie déterminée par celle-ci.

Cette profondeur n’est pas utilisée pour inférer l’hypothèse du continu. Les quatre publications
certifient l’action coinductive du même générateur ; la clôture cardinale a déjà été
obtenue au moyen de la couverture, de la condensation et des deux injections. Son
critère de vérification est double : compatibilité de toutes les troncatures
de \eqref{ix:eq:profundidad-cuatrirrelacion-ch} et provenance interne des
données de la roue dodécaphasique qui publie \(\alpha_{\rm HMT}\).
